\begin{tikzpicture}[
  >={Latex[length=3mm,width=2.2mm]},
  every node/.style={font=\small},
  vvec/.style={->,line width=1.2pt,vred},
  ghostline/.style={dashed,thick,ghost},
  cedge/.style={thick,surfedge}
]

% Superficie delimitata dalla curva C (a sinistra)
\fill[surrfblue,draw=surfedge,line width=1pt]
  plot [smooth cycle,tension=0.9] coordinates {
    (-1.0, 1.8) (0.3, 2.0) (1.2, 1.2) (1.4,-0.3)
    (0.9,-1.8) (-0.5,-2.1) (-1.4,-1.0) (-1.3, 0.6)
  };

% Curva C traslata (posizione futura, linea tratteggiata)
\draw[ghostline]
  plot [smooth cycle,tension=0.9] coordinates {
    (3.0, 1.6) (4.3, 1.8) (5.2, 1.0) (5.4,-0.5)
    (4.9,-2.0) (3.5,-2.3) (2.6,-1.2) (2.7, 0.4)
  };

% Linee che congiungono i due contorni (tubo di flusso) tratteggiate
\draw[ghostline] (1.2, 1.2) .. controls (2.2,1.6) .. (4.3,1.8);
\draw[ghostline] (1.4,-0.3) .. controls (2.4,0.0) .. (5.4,-0.5);
\draw[ghostline] (0.9,-1.8) .. controls (2.0,-1.9) .. (3.5,-2.3);
\draw[ghostline] (-0.5,-2.1) .. controls (1.5,-2.8) .. (3.5,-2.3);

% Interno della figura: profilo posteriore visibile (parte della curva traslata)
\draw[ghostline] plot [smooth, tension=0.8] coordinates {
  (0.4, 1.6) (1.0, 0.5) (0.9,-0.8) (0.2,-1.9)
};
% Freccia che indica il verso di percorrenza di C
\draw[surfedge,->,line width=0.8pt] (0.9,-0.3) -- (0.95,-0.55);

% Etichetta C
\node[font=\large\itshape,surfedge] at (-0.2,0.2) {C};

% Vettori v (in alto e in basso, diretti a sinistra)
\draw[vvec] (0.8,1.85) -- (-2.3,2.6) node[above left] {$\mathbf{v}$};
\draw[vvec] (0.3,-1.95) -- (-2.8,-1.2) node[below left] {$\mathbf{v}$};

\end{tikzpicture}